\subsection{窓・群決定的型付けの外にある標的}
\label{subsec:necessity}
前節の有限素性例は窓と可換群の積で扱えた．
以下の二つの標的は，局所群型付けと窓・群決定的型付け，
および窓・非可換群型付けと窓・アーベル群決定的型付けの
\emph{標的クラス}を分離する．
これらは特定の短縮手法の必要性や特徴標本長の下界ではない．
$\mathtt l$は検査を課す旗文字，$\mathtt r$は検査を
課さない旗文字である．
\[
 C_\times=\{\mathtt1^m\mathtt2^n\mathtt3^m\mathtt4^n:m,n\ge1\}.
\]
$\#_{\mathtt{ab}}(w)$を語$w$の因子$\mathtt{ab}$の出現数とする．
$\varphi_\times$は$\mathtt1,\mathtt3,\mathtt2,\mathtt4$を
それぞれ$(1,0),(-1,0),(0,1),(0,-1)\in\mathbb Z^2$へ，
ほかの文字を$(0,0)$へ送る加法的準同型とする．

\begin{lemma}[旗文字による分離]
\label{lem:flag-separation}
有限型付け$t$について$t(u)=t(v)$となる非空語
$u,v$と語$z$があり，
$\mathtt r uz,\mathtt r vz,\mathtt l uz\in L$で
$\mathtt l vz\notin L$なら，
すべての$f\ge1$について$L\notin\Cmcf{f}{t}$である．
\end{lemma}
\begin{proof}
$u,v$は$E=\mathtt r\blank_1z$を共有するが，
$F=\mathtt l\blank_1z$では受理が異なる．
アリティ1の分布同値が破れるので任意の$f$で成り立つ．
\end{proof}

\paragraph{局所群型付けの必要性．}
$\Sigma_{\mathrm A}=\{\mathtt l,\mathtt r,\mathtt a,\mathtt b,
\mathtt c,\mathtt\$, \mathtt1,\mathtt2,\mathtt3,\mathtt4\}$とし，
\[
 L_{\mathrm A}=
 \{\sigma w\mathtt\$c:
 \sigma\in\{\mathtt l,\mathtt r\},\
 w\in\{\mathtt a,\mathtt b,\mathtt c\}^+,\ c\in C_\times,\
 (\sigma=\mathtt r\ \text{または}\
 \#_{\mathtt{ab}}(w)\text{が偶数})\}.
\]
$s_{\mathrm A}=\lambda_{1,g}|_{\Sigma_{\mathrm A}^+}$とし，
$g:\Sigma_{\mathrm A}^2\to\mathbb Z_2$は
$\mathtt{ab}$を$1$，他を$0$へ写す．
すなわち$s_{\mathrm A}$は両端の文字と
$\mathtt{ab}$出現数の偶奇を記録する．

\begin{proposition}[局所群型付けと窓・群型付けを分離する標的]
\label{prop:lg-needed-target}
\begin{enumerate}[label=\textup{(\roman*)},leftmargin=*]
\item $L_{\mathrm A}\in\mathsf{MCFL}(2,2)\setminus\mathrm{CFL}$．
\item $L_{\mathrm A}\in\Cmcf{f}{s_{\mathrm A}}$（$f\ge2$），かつ
$s_{\mathrm A}(\Sigma_{\mathrm A}^+)\in\mathbf{LG}$．
\item 任意の有限型付け$t:\Sigma_{\mathrm A}^+\to T$について，
ある$k,\ell$と有限群への準同型
$g':\Sigma_{\mathrm A}^*\to A'$が存在して
$t\preceq(s_{k,\ell},g'|_{\Sigma_{\mathrm A}^+})$を満たすなら，
すべての$f\ge1$で$L_{\mathrm A}\notin\Cmcf{f}{t}$である．
特に正語像が$\mathbf{LI}\vee\mathbf{Grp}$に属する
任意の型付け$t$でこの非所属が成り立つ．
\end{enumerate}
したがって$L_{\mathrm A}$には
定理~\ref{thm:local-group-polydata}が適用できる一方，
どの型付けを選んでも
定理~\ref{thm:boundary-polydata}，
\ref{thm:window-group-polydata}，
\ref{thm:window-abelian-polydata}および
系~\ref{cor:mcfg-li-group-polydata}は適用できない．
特に$s_{\mathrm A}(\Sigma_{\mathrm A}^+)\in
\mathbf{LG}\setminus(\mathbf{LI}\vee\mathbf{Grp})$である．
\end{proposition}
\begin{proof}
(i) 正則言語
\[
 R_{\mathrm A}=
 \{\sigma w:\sigma\in\{\mathtt l,\mathtt r\},
 w\in\{\mathtt a,\mathtt b,\mathtt c\}^+,
 \ \sigma=\mathtt r\ \text{または}\
 \#_{\mathtt{ab}}(w)\text{が偶数}\}
\]
の右線形文法を$W$とする．
次のファンアウト2・ランク2の規則を追加する．
\[
\begin{aligned}
 C&\to(\xi_1^1\mathtt\$\mathtt1,\mathtt3)(W)
 \mid(\xi_1^1\mathtt1,\xi_1^2\mathtt3)(C),\\
 D&\to(\mathtt2,\mathtt4)
 \mid(\xi_1^1\mathtt2,\xi_1^2\mathtt4)(D),\\
 S&\to(\xi_1^1\xi_2^1\xi_1^2\xi_2^2)(C,D).
\end{aligned}
\]
これで$L_{\mathrm A}$を生成できる．
正則集合
$\mathtt r\mathtt a\mathtt\$\mathtt1^+\mathtt2^+
\mathtt3^+\mathtt4^+$
との交差後に先頭記号を消去すれば非CFLの$C_\times$を
得るため，$L_{\mathrm A}$も非CFLである．

(ii) $L_{\mathrm A}=Q_{\mathrm A}\cap\Pi_{\mathrm A}
\cap\varphi_\times^{-1}(0)$と書く．
$Q_{\mathrm A}=\{\mathtt l,\mathtt r\}
\{\mathtt a,\mathtt b,\mathtt c\}^+
\mathtt\$\mathtt1^+\mathtt2^+\mathtt3^+\mathtt4^+$は，
語頭・語尾も含めたstrictly 2-local言語である．
補題~\ref{lem:sl2-boundary-control}から
$(f,s_{1,1})$-タプル代入可能であり，
$s_{1,1}\preceq s_{\mathrm A}$なので
$(f,s_{\mathrm A})$でも代入可能である．
$\Pi_{\mathrm A}$を「先頭が$\mathtt r$，または先頭が
$\mathtt l$で$\#_{\mathtt{ab}}$が偶数」の言語とすると，
これは$\widehat{s}_{\mathrm A}$の型繊維の和集合である．
$\varphi_\times^{-1}(0)$との交差は
命題~\ref{prop:abelian-coset-filter}で代入可能性を保つ．
よって(ii)が従い，正語像が$\mathbf{LG}$に属することは
定理~\ref{thm:lg-characterization}から従う．

(iii) 窓幅$k,\ell$と群準同型
$g':\Sigma_{\mathrm A}^*\to A'$を任意に取る．
$K=\max\{1,k,\ell\}$，
$q$を$A'$の指数の正の偶数倍として
\[
 u=\mathtt a^K(\mathtt{abc})^q\mathtt a^K,\qquad
 v=\mathtt a^K(\mathtt{bac})^q\mathtt a^K
\]
とする．前後の境界は一致し，
$g'(u)=g'(\mathtt a)^{2K}=g'(v)$である．
しかし$\#_{\mathtt{ab}}(u)=q$は偶数で，
$\#_{\mathtt{ab}}(v)=1$は奇数である．
$z=\mathtt\$\mathtt{1234}$として
$\mathtt r uz,\mathtt r vz,\mathtt l uz\in L_{\mathrm A}$，
$\mathtt l vz\notin L_{\mathrm A}$であるから，
補題~\ref{lem:flag-separation}が適用できる．
$\mathbf{LI}\vee\mathbf{Grp}$の型付けは
補題~\ref{lem:mcfg-li-group-window}により
窓・群型付けで決定されるので，主張が従う．
\end{proof}

\paragraph{非可換群型付けの必要性．}
$\Sigma_{\mathrm B}=\{\mathtt l,\mathtt r,\mathtt a,\mathtt b,
\mathtt\$, \mathtt1,\mathtt2,\mathtt3,\mathtt4\}$，
$\alpha=(12),\beta=(23)\in S_3$とする．
$g_{\mathrm B}(\mathtt a)=\alpha$，
$g_{\mathrm B}(\mathtt b)=\beta$，
それ以外の文字には単位元を対応させる．
\[
 L_{\mathrm B}=
 \{\sigma w\mathtt\$c:
 \sigma\in\{\mathtt l,\mathtt r\},\
 w\in\{\mathtt a,\mathtt b\}^+,\ c\in C_\times,\
 (\sigma=\mathtt r\ \text{または}\ g_{\mathrm B}(w)=1)\}.
\]
$s_{\mathrm B}=(s_{1,1},g_{\mathrm B})$を正語像に制限する．

\begin{proposition}[窓・非可換群型付けと窓・可換群型付けを分離する標的]
\label{prop:nonabelian-needed-target}
\begin{enumerate}[label=\textup{(\roman*)},leftmargin=*]
\item $L_{\mathrm B}\in\mathsf{MCFL}(2,2)\setminus\mathrm{CFL}$．
\item 任意の$f\ge2$で$L_{\mathrm B}\in\Cmcf{f}{s_{\mathrm B}}$である．
$s_{\mathrm B}$は有限非可換群$S_3$を用いる
窓・群決定型付けである．
\item 任意の有限型付け$t:\Sigma_{\mathrm B}^+\to T$について，
ある$k,\ell$と有限\emph{アーベル群}への準同型
$g':\Sigma_{\mathrm B}^*\to A'$が存在し，
$t\preceq(s_{k,\ell},g'|_{\Sigma_{\mathrm B}^+})$なら，
すべての$f\ge1$で$L_{\mathrm B}\notin\Cmcf{f}{t}$である．
特に正語像が$\mathbf{LI}\vee\mathbf{Ab}$に属する
任意の$t$についても同様である．
\end{enumerate}
したがって$L_{\mathrm B}$には非可換群の定数
$c_{S_3,f}=M_{S_3,f}-1$を用いて
定理~\ref{thm:window-group-polydata}を適用できるが，
定理~\ref{thm:boundary-polydata}と
\ref{thm:window-abelian-polydata}は，
どの型付けの選択でも適用できない．
\end{proposition}
\begin{proof}
(i) $S_3$の値を有限状態で追う右線形文法を
上の$W$の代わりに使うと，同じMCFGで生成できる．
同じ正則交差と準同型で非CFL性を得る．
(ii) 構造を指定するstrictly 2-local言語，
旗文字と$g_{\mathrm B}$の型繊維の和集合，
および$\varphi_\times^{-1}(0)$の交差として表せる．
補題~\ref{lem:sl2-boundary-control}と
命題~\ref{prop:abelian-coset-filter}から代入可能性が従う．
(iii) 任意の可換群準同型と窓幅$k,\ell$を取り，
$K\ge\max\{1,k,\ell\}$を偶数として
\[
 u=\mathtt a^K(\mathtt{ab})^3\mathtt a^K,\qquad
 v=\mathtt a^K(\mathtt{ab})^2(\mathtt{ba})\mathtt a^K
\]
を取る．両語はParikhベクトルと前後境界が一致するので
任意の窓・可換群型付けでは同じ型となる．しかし
$p=\alpha\beta$は位数3，$\beta\alpha=p^{-1}$なので
\[
 g_{\mathrm B}(u)=p^3=1,\qquad
 g_{\mathrm B}(v)=p^2p^{-1}=p\ne1.
\]
$z=\mathtt\$\mathtt{1234}$とした
旗文字分離補題で非所属を得る．
$\mathbf{LI}\vee\mathbf{Ab}$の型付けも
窓・可換群型付けで決定される．
\end{proof}
\begin{remark}[二つの分離例が示すこと]
\label{rem:necessity-scope}
命題~\ref{prop:lg-needed-target}と
\ref{prop:nonabelian-needed-target}が分離するのは
\emph{型付けクラス}である．
局所群短縮や非可換群短縮が，
これらの標的の多項式特徴データに不可欠とは主張しない．
実際，正則言語$Q_{\mathrm A}$の構文的準同型の正語制限を
$\sigma_{\mathrm Q}$とし，正語像上で
$s'=(s_{\mathrm A},\sigma_{\mathrm Q})$と置く．
$Q_{\mathrm A}\cap\Pi_{\mathrm A}$は$\widehat{s'}$で認識されるため，
ある受理集合$\mathsf{Acc}'$について
\[
 L_{\mathrm A}
 =\widehat{s'}^{-1}(\mathsf{Acc}')\cap\varphi_\times^{-1}(0)
\]
と書ける．$s_{\mathrm A}\preceq s'$なので
$L_{\mathrm A}\in\Cmcf{f}{s'}$であり，
命題~\ref{prop:lg-needed-target}(i)の文法と
定理~\ref{thm:mcfg-abelian-filter-data}により，
$s'$の代数構造に仮定を課さなくても
$\mathcal B^{\le2}_{f,s'}$の多項式特徴データが得られる．
$L_{\mathrm B}$でも同様である．
しかも分離はアリティ1で起こる
（補題~\ref{lem:flag-separation}）．
交差依存部$C_\times$は標的を非CFLにする役割を担うが，
型付け依存部は単一成分にある．
両例は特徴標本長の下界を与えない．
\end{remark}
